\documentclass{article}

\usepackage{amsmath}

\title{Ch. 25 CW Solutions}
\author{Principles of Macroecomonics}
\date{\today}

\begin{document}

\maketitle

\begin{enumerate}

\item
	\[ \$600 \left( \frac{300}{200} \right) = \$900 \]
	
\item Nominal interest rate: \[ \left( \frac{\$2,100 - \$2,000}{\$2,000} \right) 100 = 5\% \]
	Inflation rate: \[ \left( \frac{204 - 200}{200} \right) 100 = 2\% \]
	Real interest rate: 5\% - 2\% = 3\%

\item 

	\begin{enumerate}
	
	\item Cost of Basket:
		\begin{equation*}
		12(25(\$11) + (\$700) + 1(\$120) + 1(\$40)) = 12(\$1,135)
		\end{equation*}
		Cost of Basket in Base Year: 
		\begin{equation*}
		12(25(\$10) + (\$600) + 1(\$100) + 1(\$50)) = 12(\$1,000)
		\end{equation*}
		CPI:
		\begin{equation*}
		 \frac{12(\$1,135)}{12(\$1,000)} \cdot 100 = \frac{\$1,135}{\$1,000} \cdot 100 = 113.5
		\end{equation*}
		
	\item \begin{equation*}
		\left( \frac{113.5 - 100}{100} \right)100 = 113.5\%
		\end{equation*}
		
	\item They are worse off because their nominal income rose less than the price level. That means their purchasing power decreased.
		
	\end{enumerate}
	
\item Real income in 1985:
	\begin{equation*}
	\$23,618 \left( \frac{100}{107.6} \right) = \$21,949.81
	\end{equation*}
	Real income in 1995:
	\begin{equation*}
	\$34,075 \left( \frac{100}{152.4} \right) = \$22,358.92
	\end{equation*}
	Real income in 2005:
	\begin{equation*}
	\$46,326 \left( \frac{100}{195.3} \right) = \$23,720.43
	\end{equation*}
	Real income in 2015:
	\begin{equation*}
	\$57,230 \left( \frac{100}{237} \right) = \$24,147.68
	\end{equation*}
	
	The median family's purchasing power increases in every decade.
	
\item

	\begin{enumerate}
	
	\item
		Cost of basket in 2022:
		\begin{equation*}
		30(\$4) + 10(\$8) + 10(\$8) = \$280
		\end{equation*}
		Cost of basket in 2023:
		\begin{equation*}
		30(\$4) + 10(\$9) + 10(\$10) = \$310
		\end{equation*}
		Percent change:
		\begin{equation*}
		\frac{310 - 280}{280} \cdot 100 = 10.7\%
		\end{equation*}
		
	\item In 2023, two chickens cost less than one ham. As a result, consumers will buy zero hams and twenty additional chickens. Therefore, the cost of their actual consumption basket is
	\begin{equation*}
	50(\$4) + 0(\$9) + 10(\$10) = \$300.
	\end{equation*}
	That means the actual percent change in the cost of consumption is
	\begin{equation*}
	\frac{300 - 280}{280} \cdot 100 = 7.1\%,
	\end{equation*}
	so substitution bias made my answer to (a) 3.6 percentage points too high.
	
	\end{enumerate}
	
\item In 2025, Albert has
	\begin{equation*}
	\$1,000(1.06)^3 = \$1,191.02.
	\end{equation*}
	In real dollars, that's
	\begin{equation*}
	\$1,191.02 \left( \frac{100}{118} \right) = \$1,009.34.
	\end{equation*}
	Therefore, his purchasing power increased
	\begin{equation*}
	\frac{\$1,009.34 - \$1,000}{\$1,000} \cdot 100 = 1\%.
	\end{equation*}
	
\item They expect inflation in year one to be
	\begin{equation*}
	\frac{110 - 100}{100} \cdot 100 = 10\%,
	\end{equation*}
	and they expect inflation in year two to be
	\begin{equation*}
	\frac{121 - 110}{110} \cdot 100 = 10\%.
	\end{equation*}
	For Frank to get a real annual interest rate of 2\%, he needs to charge a nominal annual interest rate of
	\begin{equation*}
	2 + 10 = 12\%.
	\end{equation*}
	
\item Let's say that this year is our base year, so today's CPI is 100. Then, the CPI in 16 years is $100(1.06)^16 = 254.04$. That means her one-time payment will be worth
	\begin{equation*}
	\$20,000 \left( \frac{100}{254.04} \right) = \$7,872.93
	\end{equation*} 
	
\item The real value of a home in 1963 was
	\begin{equation*}
	\$18,000 \left( \frac{100}{30.6} \right) = \$58,823.53,
	\end{equation*}
	and in 2025, it was
	\begin{equation*}
	\$426,800 \left( \frac{100}{322} \right) = \$132,546.60.
	\end{equation*}
	Therefore, the average real rate of return on homeownership was
	\begin{equation*}
	\left( \left( \frac{\$132,546.60}{\$58,823.53} \right)^{\frac{1}{2025 - 1963}} - 1 \right) \cdot 100 = 1.32\%.
	\end{equation*}
	
	It appears that the average real rate of return on a house is lower than long-term Treasury bonds. 
	


\end{enumerate}

\end{document}